\chapter{Open Questions}
\label{ch:questions}

The questions below are unresolved architectural decisions. Each entry states \textbf{(a)} why the
question matters for the design, \textbf{(b)} the options under consideration, and \textbf{(c)} what
additional information is needed to close it. Decisions must be recorded in \texttt{.squad/decisions.md}
before core contract freeze.

\begin{enumerate}[label=\textbf{Q\arabic*.}]

  \item \textbf{Can the runtime execute v1 runbooks without prior migration?}

    \textit{Why it matters:} Existing gert users have production runbooks in \texttt{runbook/v0} and
    \texttt{runbook/v1} format. If the runtime can only execute current-schema runbooks, adoption requires a
    migration step that may be blocked by organizational change control. A compatibility execution mode
    (``v1 shim'') removes this barrier but adds permanent maintenance cost to the parser.

    \textit{Options:} (a)~Full compatibility mode: runtime transparently executes v1 runbooks by
    internally upgrading them to schema at load time, with no schema file written. (b)~Migration
    gate:, refuses to run v1 runbooks and provides a clear error pointing to \texttt{gert migrate}.
    (c)~Opt-in compatibility flag: \texttt{--compat-legacy} enables the shim; production runs the new schema only.

    \textit{Information needed:} A survey of how many v1-format runbooks exist in representative
    deployments, and whether \texttt{gert migrate} can be made fully lossless for all v1 constructs.

  \item \textbf{What is the concurrency model for the core runtime?}

    \textit{Why it matters:} The v1 runtime is single-run-per-process. Whether supports multiple
    concurrent runs in a single process affects the state isolation model, the trace writer design,
    the governance enforcer (is it per-run or global?), and whether \texttt{gert serve} can multiplex
    concurrent VS~Code sessions.

    \textit{Options:} (a)~Single-run-per-process, matching v1 semantics; each \texttt{gert exec} or
    \texttt{gert tui} is a separate OS process. (b)~Multiple concurrent runs in-process, with
    run-scoped state isolation via Go context and per-run goroutine groups. (c)~Multiple runs via a
    long-lived daemon (\texttt{gert serve} as the host), with runs as logical sessions.

    \textit{Information needed:} Whether the VS~Code extension use case requires concurrent runs in a
    single \texttt{gert serve} process (e.g., multiple open runbook panels), and what the goroutine
    safety requirements are for the trace writer and governance layer.

  \item \textbf{Is the trace format backward-compatible with the v1 JSONL trace, or is it a new format?}

    \textit{Why it matters:} Operators and compliance tools may have parsers, SIEM integrations, or
    dashboards built against the v1 JSONL event schema. A breaking change in the trace format requires
    those integrations to be updated before, can be rolled out. A backward-compatible superset format
    avoids this but constrains the event schema design.

    \textit{Options:} (a)~New format:, defines a clean JSONL schema with a required
    \texttt{schema\_version} field; v1 trace parsers are unsupported. (b)~Backward-compatible superset:
    trace events include all v1 fields plus new fields; existing parsers continue to work on the
    subset they understand. (c)~Parallel trace output:, writes both a trace and an optional v1
    compatibility trace (flag-enabled).

    \textit{Information needed:} An inventory of known downstream trace consumers and whether the v1
    format has fields that would be structurally incompatible with the event model.

  \item \textbf{How do extension-contributed policies compose with host-level governance policies?}

    \textit{Why it matters:} The host enforces a governance policy (command allowlist, env blocking,
    redaction). An extension may also wish to contribute policies---e.g., a security extension that
    adds additional command denylists or redaction patterns. The composition model determines whether
    extensions can weaken host policies (prohibited), extend them (likely desired), or operate in a
    separate policy namespace (safest but most limited).

    \textit{Options:} (a)~Additive composition only: extensions may only add restrictions (denylist
    entries, redaction patterns); they may never remove host-level restrictions. (b)~Policy
    namespacing: each extension has an isolated policy scope; the host applies all active policies
    via intersection. (c)~Policy delegation: the host can explicitly delegate a policy domain to an
    extension (e.g., redaction rules owned by a data-classification extension).

    \textit{Information needed:} A threat model analysis of what happens when a malicious or
    compromised extension contributes a policy that attempts to weaken host restrictions. This should
    feed into \S\ref{ch:security}.

  \item \textbf{Does the \texttt{.provider.yaml} input provider framework survive unchanged?}

    \textit{Why it matters:} v1's input provider model (\texttt{.provider.yaml}, JSON-RPC resolution,
    \texttt{from:} bindings, workspace config) is a complete and working framework. It must either be
    carried forward as-is, redesigned to fit the component model, or replaced by an extension-based
    resolution mechanism. The choice affects the schema (\S\ref{ch:schema}), the extension runtime
    (\S\ref{ch:extension}), and any operator who has a custom provider binary.

    \textit{Options:} (a)~Preserve as-is: \texttt{.provider.yaml} format and JSON-RPC protocol are
    unchanged; providers are a stable first-class concept. (b)~Redesign as extension capability:
    providers become a named capability that extensions may implement; the \texttt{.provider.yaml}
    format is deprecated. (c)~Hybrid: \texttt{.provider.yaml} providers continue to work as a
    built-in extension type, with a migration path to the full extension protocol.

    \textit{Information needed:} Whether any v1 provider binaries are in active production use that
    would be broken by a protocol change, and whether the extension handshake protocol can be made a
    strict superset of the provider JSON-RPC protocol.

  \item \textbf{Is the \texttt{gert serve} JSON-RPC contract a versioned breaking change or an
    evolutionary extension?}

    \textit{Why it matters:} The VS~Code extension consumes the \texttt{gert serve} JSON-RPC~2.0 API
    over stdio. If, changes the method names, parameter shapes, or notification payloads, the
    extension must be updated simultaneously. A versioned negotiation protocol (client declares
    supported API version at handshake) allows the extension to evolve independently of the binary.

    \textit{Options:} (a)~Hard break:, defines a new JSON-RPC method set; the VS~Code extension is
    updated as part of the launch. (b)~Versioned negotiation: the \texttt{initialize} handshake
    includes a protocol version field; the host advertises capabilities and the client selects a
    compatible version. (c)~No change: the v1 JSON-RPC surface is preserved unchanged in the schema.

    \textit{Information needed:} A complete method inventory of the v1 JSON-RPC API surface (all
    methods and notification shapes currently used by the VS~Code extension), and a determination of
    which methods require changes to expose new features (e.g., saga events, OTel trace IDs).

  \item \textbf{What is the extension supply chain and signing model?}

    \textit{Why it matters:} Extensions run in separate processes but are granted capabilities that
    include tool registration, schema field contribution, and event subscription. A malicious or
    tampered extension binary could exfiltrate captured secrets, inject commands, or suppress
    governance events. An extension signing model (verified publisher key) raises the bar for supply
    chain attacks, but adds operational friction for local development.

    \textit{Options:} (a)~No signing required; extensions are trusted if they are present in the
    declared extension directory. (b)~Optional signing: extensions may carry a detached signature;
    the host warns but does not block unsigned extensions. (c)~Mandatory signing for capability
    classes above a threshold (e.g., extensions that request \texttt{governance.policy} capability
    must be signed). (d)~Keyless signing via Sigstore/cosign for developer convenience.

    \textit{Information needed:} A threat model for the extension trust boundary (see Q4 above), and
    a survey of whether the target operator community has PKI infrastructure that could issue extension
    signing keys.

  \item \textbf{Is a gRPC transport for the VS~Code extension viable for v0, or should it be deferred?}

    \textit{Why it matters:} The current JSON-RPC~2.0 over stdio transport works but limits the
    VS~Code extension to a single active connection and requires the extension to manage the gert
    process lifecycle. A gRPC transport would enable streaming (server-push events without polling),
    multiple concurrent clients, and connection multiplexing. However, it adds protocol complexity
    and requires a port allocation model.

    \textit{Options:} (a)~Retain JSON-RPC stdio as the sole transport; gRPC deferred to a
    future version. (b)~Add gRPC as an opt-in alternative transport alongside JSON-RPC stdio.
    (c)~Replace JSON-RPC stdio with gRPC as the primary VS~Code transport; provide a backward-compat
    shim for the transition period.

    \textit{Information needed:} Benchmarking of event delivery latency over JSON-RPC stdio versus
    gRPC under realistic VS~Code workloads (step execution at 100ms/step, TUI refresh at 60fps
    equivalent). Also requires a determination of whether VS~Code's Node.js extension host supports
    gRPC client libraries without native module issues.

  \item \textbf{What is the parallel execution model for independent steps within a single run?}

    \textit{Why it matters:} v1 executes all steps sequentially. Many real runbooks contain
    independent diagnostic steps that could execute in parallel, reducing total wall-clock time.
    However, parallelism complicates the trace ordering guarantee, the governance enforcer (can two
    parallel steps both consume the same approval token?), and the TUI rendering model.

    \textit{Options:} (a)~No parallelism for now; fan-out deferred to a future version. (b)~Explicit
    parallel blocks via a new \texttt{parallel:} step type that joins on completion of all children.
    (c)~Implicit parallelism inferred from the step dependency graph (steps with no data dependency
    on preceding steps execute concurrently).

    \textit{Information needed:} Whether any v1 runbooks in active use have independent steps that
    would benefit from parallelism, and what the governance implications are for concurrent step
    execution (e.g., two CLI steps executing simultaneously under the same allowlist).

  \item \textbf{How should the saga compensation registry interact with the trace and approval model?}

    \textit{Why it matters:} A compensation registered for a step must itself be subject to
    governance (it may execute privileged commands), must appear in the trace (compensations are
    audit events), and may require its own approval gate if the original step had one. The
    interaction between the compensation executor and the governance layer is non-trivial and must
    be specified before implementation.

    \textit{Options:} (a)~Compensations execute with the same governance context as the original
    step (same allowlist, same approver roles). (b)~Compensations are exempt from approval gates
    (compensations are emergency rollback actions; blocking them on approval would defeat their
    purpose). (c)~Compensation governance is separately configurable per compensation definition.

    \textit{Information needed:} Industry precedent from Temporal (compensation activities run in
    the same workflow context) and AWS Step Functions~\cite{aws-step-functions} (catch/retry blocks have no separate IAM
    policy). A decision here feeds into \S\ref{ch:schema} (compensation field design) and
    \S\ref{ch:security} (governance scope).

  \item \textbf{What is the schema namespace convention for extension-contributed fields, and how are
    they validated?}

    \textit{Why it matters:} The schema design (\S\ref{ch:schema}) states that extension fields
    are namespaced, but the namespace convention and validation mechanism are undefined. If two
    extensions contribute fields under the same namespace, or if the host does not validate
    extension-contributed fields against a schema, runbooks using those fields will have
    unpredictable behavior at runtime.

    \textit{Options:} (a)~Extension ID as namespace prefix (\texttt{x-myext.fieldName}); validated
    against a JSON Schema fragment published by the extension at registration time. (b)~Dot-notation
    namespace under a reserved top-level key (\texttt{extensions.myext.fieldName}); host validates
    at parse time. (c)~Extension fields are opaque to the host; the extension validates its own
    fields at execution time.

    \textit{Information needed:} Whether JSON Schema Draft~2020-12 supports dynamic
    \texttt{\$ref} resolution at validation time, which would allow the host to assemble a composite
    schema from the host schema plus all registered extension schema fragments.

  \item \textbf{What is the timeout and escalation model for human approval steps?}

    \textit{Why it matters:} v1 approval gates block indefinitely waiting for a role-holder to
    approve. In production, a stalled approval can block an incident response indefinitely. SLA
    enforcement---timeout after N minutes, escalate to the next role, optionally auto-approve or
    auto-reject---is a standard pattern in ITIL emergency change processes and is implemented by
    PagerDuty~\cite{pagerduty-automation} and AWS Step Functions~\cite{aws-step-functions} human approval tasks. Without it, gert is not usable for SLA-bound
    incident response (\textit{Dennis Research Brief \S5.2}).

    \textit{Options:} (a)~No timeout; approval escalation deferred. (b)~Optional
    \texttt{timeout} field on approval steps; on expiry, run is suspended and operator is notified.
    (c)~Full escalation chain: \texttt{timeout} triggers escalation to a secondary role list; if that
    also times out, the step enters a configurable terminal state (auto-approve, auto-reject, or
    suspend-for-manual-review).

    \textit{Information needed:} Whether the runtime concurrency model (Q2 above) supports
    timer-based step suspension without blocking the process, and what notification mechanism is
    available to alert escalation targets (stdout, webhook, or an extension-contributed notifier).

\end{enumerate}
